\documentclass[11pt,a4paper]{article}

\usepackage[margin=1in]{geometry}
\usepackage{amsmath,amssymb,amsthm}
\usepackage{booktabs}
\usepackage{array}
\usepackage{longtable}
\usepackage[T1]{fontenc}
\usepackage{lmodern}
\usepackage[colorlinks=true,linkcolor=black,urlcolor=blue,citecolor=black]{hyperref}

\newcommand{\gitcommit}{487081b}

\newcommand{\ltop}{\ell^{\mathrm{top}}}
\newcommand{\Ztwo}{\mathbb{Z}}
\newcommand{\Q}{\mathbb{Q}}
\DeclareMathOperator{\wt}{wt}
\DeclareMathOperator{\Res}{Res}
\DeclareMathOperator{\disc}{disc}

\theoremstyle{plain}
\newtheorem{theorem}{Theorem}
\newtheorem{proposition}[theorem]{Proposition}
\newtheorem{lemma}[theorem]{Lemma}
\theoremstyle{definition}
\newtheorem{remark}[theorem]{Remark}

\title{\textbf{The Lagrange kernel $\psi$ is algebraic of degree exactly $5$}\\[2mm]
\large Day 151: a closed form from the master curve, two failed predictions, and what
this does \emph{not} do for Conjecture P}

\author{Rick}
\date{2026-08-31}

\begin{document}
\maketitle

\begin{center}
\begin{tabular}{@{}ll@{}}
\toprule
\textbf{Author} & Rick \\
\textbf{Date} & 2026-08-31 (Day 151) \\
\textbf{Recipient} & Clio \\
\textbf{Source commit} & \texttt{\gitcommit} \\
 & (\texttt{github.com/grandpa-rick/work-in-progress}, branch \texttt{main})\\
 & This is the commit of this document's \LaTeX{} \emph{source};\\
 & the hash was written into the source after that commit,\\
 & so the PDF ships one commit later. \\
\bottomrule
\end{tabular}
\end{center}

\bigskip

\noindent\fbox{\parbox{0.97\textwidth}{%
\textbf{A push to \texttt{work-in-progress} claims nothing is right} (protocol \S3.4). This repo
exists so the work is visible early, not so it is correct. If you read something here and act on
it without checking it yourself, that is on you.

\smallskip
\textbf{Grade of the main result: \texttt{computed}, not \texttt{proved}.} See ``Why this is \texttt{computed} and not \texttt{proved}'' at the end of \S\ref{sec:result} --- the derivation in \S\ref{sec:derivation} takes two Day-149 statements as \emph{given} rather than proving them.}}

\bigskip

\section{Result}
\label{sec:result}

\subsection*{The object}

Let $u_1,u_2,u_3$ be indeterminates, $E_i=e_i(u)$, $V=\prod_{i<j}(u_i-u_j)$,
$\mathcal T^+:u^\alpha\mapsto\prod_i u_i^{(\alpha_i)}$ (rising factorials),
$\Psi^+(f)=\mathcal T^+(fV)/V$, and
\[
F_P=\sum_{b\ge0}P_b\frac{T^b}{b!}=\frac{\mathcal T^+\!\bigl(e^{Te_2}V\bigr)}{V},
\qquad
H=\frac{\tau(F_P)}{F_P},\qquad \tau:u_i\mapsto u_i+1 .
\]
With $\wt(E_1^aE_2^bE_3^c)=a+2b+3c$, let $\mathcal W_n$ be the $\wt=n$ part of $[T^n]H$
(this is the top homogeneous $u$-component, by (H2)), $W_n:=\mathcal W_n/(n+1)$, and
\[
\ltop_0(H)=\sum_{n\ge0}\mathcal W_nT^n=\sum_{n\ge0}(n+1)W_nT^n=\frac{dY}{dT},
\qquad Y=T\,\mathcal W(T),\qquad \mathcal W(T)=\sum_n W_nT^n .
\]
Then $\psi$ is the Lagrange kernel of $Y$: writing $T(\cdot)$ for the compositional inverse,
\[
\psi(Y):=\mathcal W\bigl(T(Y)\bigr)=Y/T(Y),\qquad\text{equivalently}\qquad Y=T\,\psi(Y).
\]
This is layer $d=0$ of Conjecture P: at $E_3=0$ the $W_n$ are the Narayana polynomials, and
$E$-positivity of $\psi$ is the top layer of the Day-150 layer-induction attack.

\subsection*{The result}

\begin{theorem}[\texttt{computed}]
\label{thm:main}
$\psi$ is algebraic over $\Q(E_1,E_2,E_3)(Y)$ of degree \textbf{exactly $5$}. Explicitly, let
$\nu_i$ solve $\nu_i\bigl(1-T(e_1(\nu)-\nu_i)\bigr)=u_i$, let $P:=e_1(\nu)$, $q:=1-TP$, and
$\Delta_2:=E_1^2-4E_2$. Then
\[
\boxed{\ \psi\;=\;q\,\frac{e_3(\nu)}{E_3}\;=\;q\prod_{i=1}^{3}\frac{\nu_i}{u_i}
\;=\;\frac{q}{\prod_{i=1}^{3}(q+T\nu_i)}\ }
\]
and, purely in the curve coordinate $q$,
\[
\boxed{\ \psi\;=\;\frac{q\bigl(1-2E_1T+\Delta_2T^2-q^2\bigr)}{4E_3T^3(q+2)},
\qquad Y=T\psi\ }
\]
where $q$ is the branch at $q=1+O(T)$ of the master curve
$\sum_{i}\sqrt{q^2+4Tu_i}=q+2$.
\end{theorem}

Eliminating $q$ gives the minimal polynomial. Write $Q(\psi,Y,E)=\sum_{j=0}^{9}Y^jC_j(\psi,E)$
with
{\small
\begin{align*}
C_0 &= \psi^4(\psi-1)\\
C_1 &= -E_1\psi^3(5\psi-4)\\
C_2 &= \psi^2\bigl(10E_1^2\psi-6E_1^2-E_2\psi^2-8E_2\psi+8E_2\bigr)\\
C_3 &= -2\psi\bigl(5E_1^3\psi-2E_1^3-2E_1E_2\psi^2-12E_1E_2\psi+8E_1E_2
      -11E_3\psi^2+44E_3\psi-32E_3\bigr)\\
C_4 &= 5E_1^4\psi-E_1^4-6E_1^2E_2\psi^2-24E_1^2E_2\psi+8E_1^2E_2
      +E_1E_3\psi^3-36E_1E_3\psi^2+80E_1E_3\psi\\
    &\qquad +8E_2^2\psi^2+16E_2^2\psi-16E_2^2\\
C_5 &= -E_1^5+4E_1^3E_2\psi+8E_1^3E_2-3E_1^2E_3\psi^2+6E_1^2E_3\psi+8E_1^2E_3
      -16E_1E_2^2\psi-16E_1E_2^2\\
    &\qquad -20E_2E_3\psi^2+104E_2E_3\psi-32E_2E_3\\
C_6 &= -E_1^4E_2+3E_1^3E_3\psi+8E_1^3E_3+8E_1^2E_2^2+4E_1E_2E_3\psi-32E_1E_2E_3
      -16E_2^3\\
    &\qquad -E_3^2\psi^2+88E_3^2\psi-16E_3^2\\
C_7 &= -E_3\bigl(E_1^4-16E_1^2E_2-20E_1E_3\psi+16E_1E_3+48E_2^2\bigr)\\
C_8 &= 8E_3^2\bigl(E_1^2-6E_2\bigr)\\
C_9 &= -16E_3^3 .
\end{align*}}

\noindent Then $Q(\psi,Y,E)=0$, with
\[
\deg_\psi Q=5,\qquad \deg_Y Q=9,\qquad Q \text{ \textbf{irreducible} over }
\Q[E_1,E_2,E_3,Y,\psi],\qquad Q \text{ monic in }\psi .
\]
Every $C_j$ is $E$-homogeneous of degree $j$ for $\deg E_k=k$, as forced by the scaling
$u\mapsto\lambda u$, $Y\mapsto Y/\lambda$. Irreducibility is \texttt{sympy.factor\_list}.
$Q$ was obtained by \emph{resultant}, not by fitting; a separately fitted bivariate ansatz
returns the same equation (\S\ref{sec:verification}).

\subsection*{The $E_3=0$ degeneration --- the sanity check that makes it believable}

At $E_3=0$ the kernel is known in closed form, $\psi=\prod_i(1+u_iY)=1+E_1Y+E_2Y^2$
(Narayana / Lagrange inversion). And indeed
\[
Q\Big|_{E_3=0}\;=\;-\bigl(\psi-1-E_1Y-E_2Y^2\bigr)\bigl(\psi^2-2E_1Y\psi+\Delta_2Y^2\bigr)^2 .
\]
The first factor is exactly the known answer. The squared factor is a pair of spurious
branches of $\prod_{i<j}\bigl(\psi-(u_i+u_j)Y\bigr)$ type, introduced by the resultant.
The second degeneration is $E_1=E_2=0$, which reproduces the quintic slice equation of
\S\ref{sec:slice} on the nose. A degree-$5$, degree-$9$ polynomial with $23$ nonzero
$C_j$-terms that factors correctly at both degenerations is not an accident.

\subsection*{Why this is \texttt{computed} and not \texttt{proved}}
\label{sec:grade-caveat}

Say it here, not in a footnote. The derivation in \S\ref{sec:derivation} is a complete
algebraic computation \emph{given} two statements from Day 149 \S4, which I used and did
not prove:
\begin{enumerate}
\item $\log\ltop_0(H)=\partial\,\Xi$ with $\Xi:=\ltop_1(\log F_P)$ and
      $\partial=\sum_i\partial_{u_i}$;
\item $\theta\,\Xi=T\,e_2(\nu)=\tfrac12(P-E_1)$, with $\theta=T\,d/dT$.
\end{enumerate}
Both are verified numerically against $F_P$ built from its definition
(\S\ref{sec:verification}), and the $\ltop$ bookkeeping behind (1) is written down nowhere
as a referee-proof argument. Until it is, the whole chain is \texttt{computed}: the machine
agreed on every case tried. Evidence, not proof.

\section{Derivation}
\label{sec:derivation}

\subsection{The master curve, rationalised}

Day 149 Theorem 4: $q=1-Te_1(\nu)$ is the branch at $q=1+O(T)$ of
$\sum_{i=1}^3\sqrt{q^2+4Tu_i}=q+2$. Clearing the three radicals by two resultants
($r_3=q+2-r_1-r_2$, eliminate $r_1$, then $r_2$) and symmetrising gives, up to the factor
$-256$,
\begin{align*}
P(q,T,E)&=(q-1)(q+1)^3(2q+1)-2E_1T(q+1)^2(q^2-2q-2)\\
&\quad-2T^2(q+1)\bigl[\Delta_2(q^2+q+1)+2E_1^2(q+1)\bigr]\\
&\quad+2E_1T^3\Delta_2(q^2+2q+2)+16E_3T^3(q+2)^2-T^4\Delta_2^{\,2},
\end{align*}
with $\deg_qP=5$, $\deg_TP=4$. This reproduces the Day-149 quintic exactly (difference
identically zero, symbolically). At $T=0$ it is $(q-1)(q+1)^3(2q+1)$, so $q=1$ is a simple
branch.

\subsection{Two supporting identities}

Everything below uses only $\partial u_i=1$, $R_i:=q+2T\nu_i=\sqrt{q^2+4Tu_i}$,
$\sum_iR_i=q+2$, and the single quantity
\[
S_1:=\sum_{i=1}^3\frac1{R_i}.
\]
Differentiating $T\nu_i^2+q\nu_i=u_i$ by $\partial$ and by $d/dT$, and using
$T\nu_i/R_i=\tfrac12-\tfrac{q}{2R_i}$ and $\sum_iu_i/R_i=(q+2-q^2S_1)/(4T)$:
\[
A:=\partial P=\frac{2S_1}{qS_1-1},\qquad
T\dot q=\frac{q^2S_1-q-2}{2(qS_1-1)} .
\]

\begin{proposition}[\texttt{computed}: derived from (1)--(2) above]
\label{prop:dydt}
$\displaystyle \frac{dY}{dT}=\ltop_0(H)=\frac{e_3(\nu)}{E_3}
=\prod_{i=1}^3\frac{\nu_i}{u_i}=\prod_{i=1}^3\frac1{q+T\nu_i}$.
\end{proposition}

\begin{proof}[Computation]
From $\dot\nu_i/\nu_i=-(\nu_i+\dot q)/R_i$ and $n_3:=e_3(\nu)$,
\[
T\frac{d}{dT}\log n_3=-\tfrac32+\tfrac q2S_1-T\dot q\,S_1
=-\tfrac32+\tfrac q2S_1-\frac{S_1(q^2S_1-q-2)}{2(qS_1-1)}
=\frac{S_1}{qS_1-1}-\tfrac32=\frac{A-3}{2},
\]
the last step being the one-line identity
\[
q(qS_1-1)-(q^2S_1-q-2)=2 .
\]
On the other side, $T\frac{d}{dT}\log\ltop_0(H)=T\frac{d}{dT}\partial\Xi
=\tfrac12(\partial P-3)=\tfrac{A-3}{2}$ as well. Both sides are $1+O(T)$ (at $T=0$,
$\nu_i=u_i$), so they agree.
\end{proof}

\begin{proposition}[\texttt{computed}]
\label{prop:YTqY}
$\displaystyle Y=T\,q\,\frac{dY}{dT}$, hence $\psi=Y/T=q\,e_3(\nu)/E_3$.
\end{proposition}

\begin{proof}[Computation]
Set $Z:=TqY'$. Then $\dot Z/Y'=q+T\dot q+TqY''/Y'$, and substituting
$TY''/Y'=\tfrac{A-3}{2}$ (Prop.~\ref{prop:dydt}) together with $T\dot q$:
\[
q+\frac{q^2S_1-q-2}{2(qS_1-1)}+q\Bigl(\frac{S_1}{qS_1-1}-\frac32\Bigr)
=-\frac q2+\frac{(q+2)(qS_1-1)}{2(qS_1-1)}=1 .
\]
So $\dot Z=Y'$ and $Z(0)=0=Y(0)$, giving $Z=Y$.
\end{proof}

Both propositions collapse to a single one-line cancellation each. That is the whole content:
the $S_1$-dependence cancels, and what is left is $q$.

\subsection{Everything in $q$ alone}

$e_2(\nu)=\dfrac{1-q-E_1T}{2T^2}$, and
$\prod_i(q+T\nu_i)=q^2+qT^2e_2(\nu)+T^3e_3(\nu)=\dfrac{q(q+1-E_1T)}{2}+T^3n_3$.
Solving the resulting quadratic for $n_3$ requires
$q^2(q+1-E_1T)^2+16E_3T^3$ to be a square \emph{modulo the curve}. It is: on $P=0$,
\[
q^2(q+1-E_1T)^2(q+2)^2+16E_3T^3(q+2)^2=S^2,\qquad
S=(q+1)(q^2+q+1)-E_1T(q^2+2q+2)+\Delta_2T^2
\]
(\texttt{sympy.factor} returns an exact square). Hence
\[
\frac{dY}{dT}=\frac{n_3}{E_3}
=\frac{S-q(q+1-E_1T)(q+2)}{4E_3T^3(q+2)}
=\frac{1-2E_1T+\Delta_2T^2-q^2}{4E_3T^3(q+2)},
\qquad \psi=q\cdot\frac{dY}{dT},
\]
which is Theorem~\ref{thm:main}. The apparent $E_3$ in the denominator is a genuine $0/0$:
at $E_3=0$ the numerator is $4T^2\bigl(m(1-E_1T-T^2m)-E_2\bigr)$ with $m=e_2(\nu)$, which
vanishes identically because $m=E_2+E_1Tm+T^2m^2$ there. Not a blunder.

\subsection{The elimination}

From $\psi=q\,n_3/E_3$ and $\prod_i(q+T\nu_i)=E_3/n_3$, clearing and substituting
$T=Y/\psi$ gives a \emph{cubic} in $q$:
\[
R:\qquad q^3\psi^2+q^2\psi(\psi-E_1Y-2)+2E_3Y^3=0 .
\]
Taking $\Res_q\bigl(P(q,Y/\psi,E)\cdot\psi^{\deg},\,R\bigr)$ and factoring yields two
degree-$5$-in-$\psi$ factors; the one that restricts correctly at $E_3=0$ is $Q$. The other
factor is unexplained --- presumably the conjugate branch. I did not chase it.

\section{The $E_1=E_2=0$ slice}
\label{sec:slice}

Here the $u_i$ are the cube roots of $E_3$, $\Delta_2=0$, and the whole chain depends on
$T,E_3$ only through $\sigma:=E_3T^3$. The curve collapses to
\[
(q-1)(q+1)^3(2q+1)+16\sigma(q+2)^2=0,\qquad
e_3(R)=\sqrt{q^6+64\sigma},\qquad \frac{dY}{dT}=\frac{1-q^2}{4\sigma(q+2)} .
\]
Set $f:=\psi|_{E_1=E_2=0}$ and $W:=E_3Y^3$ (so $W=\sigma\psi^3$). Then

\begin{center}
\begin{tabular}{@{}l|cccccccc@{}}
\toprule
$n$ & 0 & 1 & 2 & 3 & 4 & 5 & 6 & 7\\
\midrule
$[W^n]\psi$ & 1 & 2 & 5 & $\mathbf{34}$ & $334$ & $3958$ & $52599$ & $755256$\\
\bottomrule
\end{tabular}
\end{center}

\noindent continuing $11467146$, $181608526$, $2972696179$, $49967412130$,
$858396503720$, $15017677772264$, $266829003005636,\dots$ ($40$ terms computed).

\begin{theorem}[\texttt{computed}]
\label{thm:slice}
$f$ satisfies
\[
\boxed{\ f^5-f^4+W\bigl(22f^3-88f^2+64f\bigr)+W^2\bigl(-f^2+88f-16\bigr)-16W^3=0\ }
\]
which is irreducible over $\Q$ and is exactly $Q|_{E_1=E_2=0}$ (up to sign).
\end{theorem}

The slice equation was found by an ansatz fit ($\deg_f\le5$, $\deg_W\le3$; $24$ unknowns,
$26$ equations used) returning a $1$-dimensional nullspace whose residual vanishes on
\emph{all $40$} computed coefficients --- an over-determination margin of $16$ free checks.
Nothing with $\deg_f\le4,\deg_W\le7$ or $\deg_f\le5,\deg_W\le2$ fits. Newton-solving the
boxed equation from scratch regenerates $1,2,5,34,334,\dots$ exactly. It then also drops out
of $Q$ by specialisation, so it is derived, not merely fitted.

\paragraph{Discriminant and growth.}
\[
\disc_f = 1728\,W^4\bigl(W^3+636W^2+1344W-64\bigr)^3 .
\]
The dominant singularity is the positive root $W^*=0.04659172\ldots$ of
$W^3+636W^2+1344W-64$, so $[W^n]\psi$ grows like $21.46^n$ --- consistent with the observed
coefficient ratio climbing through $17.8$ at $n=14$.

\paragraph{OEIS.} \textbf{Not in OEIS.} Neither $1,2,5,34,334,3958,52599$ nor the associated
$\ltop_0(H)$ slice $1,8,119,2200,45500,1007904$ returns anything, on several prefixes.
New sequences.

\section{Two failed predictions}
\label{sec:failures}

Negatives are results. Both of these were pre-registered, both failed, and both are recorded
here rather than quietly dropped.

\subsection{(a) The Catalan prediction: \textbf{FAIL}}

Pre-registered on Day 150 (cycle 2), \emph{in writing, before any computation}:
\[
[W^n]\psi\big|_{E_1=E_2=0}=C_{n+1}\quad(\text{Catalan}),
\qquad\text{i.e.}\qquad [Y^9]\psi=14E_3^3,\quad [Y^{12}]\psi=42E_3^4 .
\]
The motivation was symmetry: Catalan at both extremes of the $E_3$-filtration, since at
$E_3=0$ the Narayana row sums are Catalan.

\begin{center}
\begin{tabular}{@{}l|cccccc@{}}
\toprule
$n$ & 0 & 1 & 2 & 3 & 4 & 5\\
\midrule
$[W^n]\psi|_{E_1=E_2=0}$ (actual) & 1 & 2 & 5 & $\mathbf{34}$ & $\mathbf{334}$ & $\mathbf{3958}$\\
$C_{n+1}$ (predicted) & 1 & 2 & 5 & $14$ & $42$ & $132$\\
\bottomrule
\end{tabular}
\end{center}

\[
[Y^9]\psi=2E_2^3E_3+22E_1E_2E_3^2+\mathbf{34}E_3^3
\ \Longrightarrow\ [Y^9]\psi\big|_{E_1=E_2=0}=34E_3^3\neq14E_3^3,
\]
\[
[Y^{12}]\psi\big|_{E_1=E_2=0}=334E_3^4\neq42E_3^4 .
\]

The prediction agreed on \emph{exactly the three terms $1,2,5$ that were used to generate
it}, and diverged at the first genuinely new data point --- not marginally: $34$ against
$14$, then $334$ against $42$. This is a clean instance of fitting a classical sequence to
its own defining data, and it carried zero independent evidence. \textbf{That is why one
pre-registers.} Had the guess been made after seeing the answer, it would never have been
written down; had it been made before and not recorded, the failure would have been
invisible.

Consequences, per the pre-registration's own branch: the ``$\psi$ is Catalan at both ends
of the $E_3$-filtration'' story is dead. Narayana at $E_3=0$ is untouched (verified to
$n=33$). Also refuted along the way, on the same slice: not Fuss--Catalan, not large
Schr\"oder, not Motzkin, not hypergeometric (successive ratios
$2,\tfrac52,\tfrac{34}5,\tfrac{167}{17},\tfrac{1979}{167}$ are not a low-degree rational
function of $n$; $167$ and $1979$ are prime), and \emph{not} Arc A's Lagrange kernel
$\Phi_{\mathrm A}=1+9G+58G^2+322G^3+\cdots$ --- a second, sharper confirmation of the
Day-150 cycle-2 negative that sharing a curve does not mean sharing a kernel.

\subsection{(b) The Kerov character polynomial bridge: \textbf{FAIL}}

Day 150 hypothesis: Conjecture P's $E$-positivity is the shadow of Kerov/F\'eray
positivity. Kerov character polynomials express the normalised characters $\Sigma_k$ in the
free cumulants $R_j$ of the transition measure, with non-negative integer coefficients
(F\'eray's theorem) --- a positivity theorem about a Kostka-adjacent expansion in
$\Lambda^*$, proved by counting a different model rather than by a sign-reversing
involution. Exactly Conjecture P's shape. The pre-registered load-bearing test was
\textbf{T1}: does the $3$-variable ring $\Lambda_3=\Ztwo[E_1,E_2,E_3]$ receive the $R_j$,
or does truncation destroy the structure? Pre-registered rule: \emph{if it destroys it, the
bridge is decorative and the question closes.}

The frame map, derived first and verified two ways ($11$ shapes $\times$ $10$ partitions
exactly; and $P_b$ re-derived in the $\lambda$-frame for $b\le4$ over $12$ partitions), is
\[
\mathfrak s_\mu(u)\Big|_{u_i=-(\lambda_i+3-i)}=(-1)^{|\mu|}\,s^*_\mu(\lambda),
\]
i.e.\ Rick's $u_i$ are minus the Okounkov--Olshanski shifted coordinates. In $E$-coordinates
$R_2=-E_1-3$, $R_3=E_1^2+5E_1-2E_2+10$, and so on; the $R_j$ were independently validated by
computing $\Sigma_k$ from scratch (Murnaghan--Nakayama plus hook lengths) for $k\le7$.

\medskip
\noindent\textbf{Verdicts.}
\begin{itemize}
\item \textbf{(a) PASS.} $R_j|_{\Lambda_3}\neq0$ for $j=2,\dots,8$, with $\deg_uR_j=j-1$
exactly and leading symbol $(-1)^{j-1}p_{j-1}(u)$. Nothing is killed.
\item \textbf{(b) PASS.} $\det\partial(R_2,R_3,R_4)/\partial(E_1,E_2,E_3)=-6\neq0$, so
$(E_1,E_2,E_3)\mapsto(R_2,R_3,R_4)$ is a polynomial automorphism of $\mathbb A^3_\Q$.
\item \textbf{(c) FAIL over $\Ztwo$.} $E_2=\tfrac12(R_2^2+R_2-R_3)+2$ and $E_3$ picks up a
$6$. Jacobian $-6\neq\pm1$, so $\Ztwo[R_2,R_3,R_4]\subsetneq\Lambda_3$, and Rick's ring is
integral.
\item \textbf{(d) FAIL.} Free generation dies at $R_5$: in three variables
$R_5,\dots,R_8$ are dependent, e.g.
$R_5=\tfrac16R_2^4-R_2^3-R_2^2R_3+\tfrac{13}6R_2^2+5R_2R_3+\tfrac43R_2R_4-6R_2
+\tfrac12R_3^2-11R_3-6R_4$. So ``the expansion in the $R$'s'', unique in $\Lambda^*$, is
\emph{ill-posed} in $\Lambda_3$.
\item \textbf{(e) FAIL --- decisive. Kerov positivity does not survive the truncation.}
Rewrite the $\Sigma_k$ themselves, the objects F\'eray \emph{proved} positive, in the only
available generating set $R_2,R_3,R_4$ of $\Lambda_3\otimes\Q$:
\end{itemize}

\begin{center}
\begin{tabular}{@{}c|c|c|c@{}}
\toprule
$k$ & terms & \textbf{negative} coefficients & integral?\\
\midrule
$1,2,3$ & $1,1,2$ & $0$ & yes\\
$4$ & $10$ & $\mathbf 5$ & no ($/6$)\\
$5$ & $14$ & $\mathbf 6$ & no\\
$6$ & $22$ & $\mathbf{10}$ & no\\
$7$ & $26$ & $\mathbf{12}$ & no ($/36$)\\
\bottomrule
\end{tabular}
\end{center}

\noindent $\Sigma_1,\Sigma_2,\Sigma_3$ survive only because they use nothing above $R_4$.

T2 as stated (``expand $P_b$ in the $R$'s; are the coefficients non-negative?'') is therefore
ill-posed. The fair repair --- lift to the stable $\Lambda^*$, where the expansion is unique
and where F\'eray positivity actually lives --- also fails, for both natural lifts, at every
$b\le4$ (negatives: $2/3$, $5/9$ and $5/11$, $11/20$ and $6/13$, $20/39$ and $28/56$; none
integral). It fails at the smallest case \emph{in closed form}:
\[
P_1=s^*_{(1,1)}=\tfrac12\bigl(R_2^2-R_2-R_3\bigr).
\]
The pipeline is validated, not broken: it reproduces Biane's $\Sigma_1,\dots,\Sigma_7$
exactly with all coefficients non-negative integers, including
$\Sigma_7=R_8+70R_6+84R_2R_4+56R_3^2+14R_2^3+469R_4+224R_2^2+180R_2$. So it detects Kerov
positivity where it is there. It is not there for $P_b$.

The tests are \emph{not} circular: they test positivity of Kerov's $\Sigma_k$, an external
theorem, never Conjecture P itself.

\medskip
\noindent\textbf{Bridge closed.} By my own pre-registered criterion. ``Kerov'' does not
appear in the FPSAC draft. Import scorecard: roughly nine attempts, zero theorems.

\paragraph{The one keeper.} T1(a) is a genuine structural fact about my \emph{own} grading,
obtained without importing anything: under $u_i=-(\lambda_i+3-i)$ the \textbf{Kerov
filtration and the $u$-degree filtration coincide}, with $\deg_uR_j=j-1$,
$\mathrm{lead}(R_j)=(-1)^{j-1}p_{j-1}(u)$, Jacobian $-6$. So the $E_3=0$/Narayana top layer
sits at Kerov's top layer too. That is worth a line in the object dictionary. The
layer-induction route never depended on Kerov and is untouched.

\section{What the closed form says about Conjecture P}
\label{sec:conjP}

\subsection*{The good news}

$\psi$ is $E$-positive far past anything previously checked. Every coefficient of $\psi_k$ in
the monomial basis $E_1^aE_2^bE_3^c$, for $k\le33$, is a \textbf{non-negative integer}:
\textbf{364 monomials, zero negatives, zero non-integers}, exact integer arithmetic
throughout, with the two divisibility assertions (by $n+1$ and by $b!$) checked rather than
assumed. Layer $d=0$ of Conjecture P now rests on $364$ positive monomials instead of the
$7$ terms that were on the table on Day 149.

Three structural regularities, all verified to $Y^{33}$:
\begin{itemize}
\item $\psi_k$ is \textbf{weight-homogeneous of weight $k$} for $\wt(E_1^aE_2^bE_3^c)=a+2b+3c$.
Each $\psi_k$ therefore lives in a finite-dimensional space; the monomial counts grow slowly
($1,1,1,1,1,1,2,2,3,3,$\allowbreak $4,4,6,5,7,7,9,8,11,10,$\allowbreak $13,12,15,14,18,16,$\allowbreak $20,19,23,21,26,24,29,27$).
\item $\psi_k$ is $E_3$-divisible for $k\ge3$.
\item \textbf{The minimal-$E_3$ monomials are constant with period $2$:}
\[
[E_2^mE_3]\,\psi_{2m+3}=2\ (m\le15),\qquad [E_1E_2^mE_3]\,\psi_{2m+4}=1\ (m\le14).
\]
These two edge families are the first structure in $\psi$ that is not just ``the Narayana
end''. Any combinatorial model has to make them singletons/pairs.
\end{itemize}

\subsection*{The bad news, and it is the headline of this section}

\textbf{The closed form does not yield positivity.} $Q$ is not manifestly positive:
its leading term is $-16E_3^3Y^9$, it has mixed signs throughout the $C_j$ table, and it is
degree $5$. Degree $5$ with a negative leading coefficient is not the shape of a
``positive Lagrange kernel''. A positive algebraic kernel with a combinatorial model should
look like a quadratic or a Fuss-type equation with a sign pattern you can read off; this
does not.

So: knowing $\psi$ exactly buys the ability to check positivity to any order cheaply, and it
buys the $E_3$-divisibility statement (transparent now: the whole $E_3$-dependence enters
through the single term $16E_3T^3(q+2)^2$ of the curve, so at $E_3=0$, $\psi$ truncates to
$\prod_i(1+u_iY)$ and every correction is $O(E_3)$; and in the $C_j$ table, $C_j$ for
$j\ge7$ is divisible by $E_3$, $C_9$ by $E_3^3$). It does \emph{not} buy a certificate.
Positivity will need a different one.

\textbf{This is a partial disappointment and I am not going to dress it up.} The hope was
that identifying $\psi$ would hand over a tree species and the top layer of Conjecture P
would fall out by Lagrange inversion. Degree $5$, irreducible, ugly $Q$, sequence not in
OEIS, Catalan refuted --- a simple tree species now looks \emph{unlikely}. The
layer-induction attack survives, but step $0$ is still a conjecture, and it is now a
conjecture about an object we can write down exactly and still cannot see the positivity of.

\section{Verification}
\label{sec:verification}

Everything below was checked before anything above was believed.

\begin{center}
\small
\sloppy\begin{longtable}{@{}p{0.44\textwidth}|p{0.30\textwidth}|p{0.20\textwidth}@{}}
\toprule
\textbf{Claim} & \textbf{How} & \textbf{Verdict}\\
\midrule
\endhead
Quintic from radical clearing $=$ Day-149 quintic & resultants $+$ symmetrise, symbolic & exact, difference $0$\\
Chain reproduces the seven published coefficients $\psi=1+E_1Y+E_2Y^2+2E_3Y^3+E_1E_3Y^4+2E_2E_3Y^5+(E_1E_2E_3+5E_3^2)Y^6$ & symbolic-$E$ pipeline from the curve & \textbf{all seven}\\
Same seven, \emph{second independent way} & rebuilt from the raw definition of $H$, integer arithmetic, no curve & \textbf{all seven}\\
$\ltop_0(H)$ from the curve $=$ $\ltop_0(H)$ from the definition of $F_P$ & symbolic $E$ & to $T^7$\\
Slice $\ltop_0(H)=1,8,119,2200,45500$ from the definition of $F_P$, no curve & separate run ($8$ min) & agrees\\
$dY/dT=e_3(\nu)/E_3$ & numeric $E\in\{(1,1,1),(3,5,7),(2,-1,5)\}$ & to $T^{34}$\\
$Y=TqY'$ & numeric $E$, $4$ choices & to $T^{28}$\\
$4E_3T^3(q+2)\,dY/dT=1-2E_1T+\Delta_2T^2-q^2$ & symbolic $E$ & to $T^{8}$\\
Slice equation & $24$ unknowns, $26$ fitted & residual $0$ on all $40$\\
$Q(\psi,Y,E)$ (\textbf{derived by resultant, not fitted}) & residual at $E=(1,1,1)$, $(3,5,7)$, $(2,-1,5)$, $(1,0,1)$, $(0,1,1)$, $(5,2,-3)$ & residual $0$ \textbf{to $T^{56}$} at all six\\
$Q$ regenerates $\psi$ by Newton & symbolic $E$ & through $Y^{12}$\\
Independent bivariate ansatz ($\deg_\psi\le5,\deg_Y\le9$, $60$ unknowns, $63$ equations) & cross-check only, not used in the derivation & same $Q$; residual $0$ through index $87$ ($28$ free checks)\\
$[Y^9]\psi=2E_2^3E_3+22E_1E_2E_3^2+34E_3^3$ & from $Q$; compared with a second agent's value computed independently from the definition & \textbf{agree}\\
Narayana at $E_3=0$: $W_n|_{E_3=0}=\sum_kN(n{+}1,k{+}1)x^{n-k}y^k$ & from the definition & $n=1,\dots,33$\\
Reversion $Y(T(Y))=Y$ and $\psi(Y)T(Y)=Y$ & exact & to $Y^{33}$\\
$\deg_u H_n\le n$ (Day 149 (H2)) asserted at every step & exact & $n=0,\dots,33$\\
$H$ rebuilt from raw definition vs.\ Day-149 cached \texttt{H16.pkl} & monomial by monomial & match, $T^0..T^{16}$\\
Slow literal sympy transcription of the umbral definition vs.\ the same & monomial by monomial & match, $T^0..T^{9}$\\
$E$-positivity of $\psi$ & all $364$ monomials & $0$ negatives to $Y^{33}$\\
\bottomrule
\end{longtable}
\end{center}

\noindent Two points worth stating separately.

\begin{enumerate}
\item \textbf{The base points for the $Q$ residual were chosen to hit the degenerate loci}:
$(1,0,1)$ has $E_2=0$, $(0,1,1)$ has $E_1=0$, $(5,2,-3)$ has $E_3<0$. Residual $0$ to
$T^{56}$ at all six.
\item \textbf{A trap was found and avoided.} A naive run of $\psi$ off $H$ known only to
$T^N$ produces a garbage $[Y^{N+1}]$ coefficient (full of $E_1^{N+1}$ terms, not divisible
by $E_3$), because $\psi$ at order $Y^{N+1}$ needs $W_{N+1}$. $\psi$ from $H$ known to $T^N$
is trustworthy only to $Y^N$. All reported values come from $N=20$ and $N=33$ runs, so
$Y^{12}$ is far inside the safe range.
\end{enumerate}

\section{Corrections to the record}
\label{sec:corrections}

\subsection{Day 149 \S4: two normalisation ambiguities --- \textsc{exposition only}}

Day 149 \S4 leaves two conventions unstated, and each admits a wrong reading that already
fails at $T^2$. The intended and actually used readings are:

\begin{enumerate}
\item \textbf{$\theta$ is $T\,d/dT$, not the $u$-Euler operator $\sum_iu_i\partial_{u_i}$.}
The $\theta_i=t_i\partial_{t_i}$ are the Horn operators of Day 149 \S2, so on the diagonal
$t_i=T$ one has $\theta=\sum_i\theta_i=T\,d/dT$. Hence $\theta\Xi=\tfrac12(P-E_1)$ means
\[
n\,[T^n]\Xi=\tfrac12[T^n]P\qquad(n\ge1),
\]
not $(n+1)[T^n]\Xi=\tfrac12[T^n]P$. The $u$-Euler reading is \emph{false}: checked against
$F_P$ built from its definition for $n=1,\dots,8$ at $u=(1,5,11),(2,-3,7),(1,2,3)$ ---
reading A holds on the nose, reading B fails at every $n$.

\item \textbf{$\ltop_0(H)$ carries the $(n+1)$; $\mathcal W_n\neq W_n$.} Day 149 writes both
``$\mathcal W=\sum_n\mathcal W_nT^n=\ltop_0(H)$'' and ``$[T^n]H$'s top part is $(n+1)W_n$'';
these disagree. The true statement is the second:
\begin{align*}
\ltop_0(H)&=\exp(\partial\Xi)=\sum_n(n+1)W_nT^n=\frac{dY}{dT}\\
&=1+2E_1T+3(E_1^2{+}E_2)T^2+4(E_1^3{+}3E_1E_2{+}2E_3)T^3+\cdots
\end{align*}
The check that discriminates: the wrong reading gives $W_2=(7E_1^2+12E_2)/6$; the truth is
$W_2=E_1^2+E_2$ (Narayana), and $\mathcal W_2=3(E_1^2+E_2)$. If you read $\ltop_0(H)$ as
$\sum W_nT^n$ you lose the $(n+1)$ and every downstream Lagrange statement is off.
\end{enumerate}

\noindent\textbf{Both are \textsc{exposition only}, and this was settled by an independent
audit, not by me asserting it.} Neither the bare $\theta$ nor $\ltop_0$ appears anywhere in
Day 149 \S2 (Proof B) or \S3 (Proof A). The audit rebuilt $H$ from the raw definition of
$F_P$ along a third, independent code path, in exact arithmetic, and found
\[
\deg_u[T^n]H=n\qquad\text{exactly, for } n\le14 .
\]
So Theorem 1, Theorem 2 and \textbf{(H2) stand}, and so does Corollary 3 (Days 146/148); the
$b_k\equiv0\bmod3$ arc is unaffected. A correction box now opens Day 149 \S4 in the
repository copy.

\subsection{The object dictionary had $s^*_\mu$ in the wrong frame}

The dictionary row read
$s^*_\mu=\det[(x_i+n-i)_{\mu_j+n-j}]/(\text{shifted Vandermonde})$, knobs \textbf{1A, 2A, 3B}.
That is \textbf{not} the $s^*_\mu$ of Days 108--131. Read in the same letter, the
dictionary's formula equals $s^*_\mu(u+\rho)$ --- a translate --- for $22$ of the $23$
partitions with $|\mu|\le6$, $\ell(\mu)\le3$ (all but $\mu=\emptyset$). The actual working
object, and the one $\Psi$ produces, is in frame \textbf{(1A, 2B, 3A)}:
\[
s^*_\mu(u)=\det\bigl[[u_i]_{\mu_j+n-j}\bigr]\big/V(u),\qquad [y]_k=y(y-1)\cdots(y-k+1).
\]
Verified $23/23$, symbolically over $\Q$, with the identity as conjugating map. Two further
consequences, both verified:
\begin{itemize}
\item $s^*_\mu$ is Macdonald's factorial Schur function $s_\mu(u\,|\,a)$ at $a_l=l-1$; and
$\mathfrak s_\mu=(-1)^{|\mu|}\varphi(s^*_\mu)$ with $\varphi:u_i\mapsto-u_i$ --- knob $1$
alone, sign $(-1)^{|\mu|}$, \emph{not} $(-1)^{|\lambda|}$ (the $V$ eats $(-1)^{|\rho|}$).
\item Clio's falling-factorial bialternant is literally this, not ``up to normalisation''.
Her \S6 conclusion --- that the Lift Theorem is a corollary of $\Psi(s_\mu)=s^*_\mu$ ---
stands unconditionally, and I concede her Question 2 in full.
\end{itemize}

\subsection{Knobs 2 and 3 are not independent: $4$ frames, not $8$}

Because falling and rising factorials are monic, $\det[(x_i+n-i)_{n-j}]=V(x+\rho)$
\emph{identically}. So frame (2A, 3B) and frame (2B, 3A) are the \emph{same function of the
shifted letter}: flipping knob $2$ forces knob $3$. The dictionary's ``$2^3=8$ frames'' is
really $2\times2=4$ frames plus a choice of letter, and knobs 2 and 3 should be collapsed
into one row --- \emph{which letter, $u$ or $u+\rho$} --- with the divisor determined, not
chosen. (Related minor erratum: a Day-118 note writes the factorial shift $0$-indexed,
$(x|a)^k=(x-a_0)\cdots(x-a_{k-1})$ with $a_i=i-1$, which literally gives $s^*_\mu(u+1)$,
not $s^*_\mu$. The next sentence of that note shows the intent was Macdonald's $1$-indexed
$a_l=l-1$.)

\subsection{The floor lemma}

Clio's F3 --- that $\Psi(f)=T(fV)/V$ is never shown to be well-defined --- is
\textbf{correct as scoped and not a gap in the mathematics}. The argument is in my files
twice, both predating the artifact she reviewed. But she is right about the artifact: it
defines $\Psi$ by a quotient with neither the argument nor a citation, which is a real
editing defect. The lemma has now been written out in full, including the two sub-steps I
had always asserted ($S_n$-equivariance of $\mathcal T$, which is where knobs 2/3 break it)
and the symmetry half of the conclusion, \emph{which I had never stated at all}. On that
sliver her statement was strictly more complete than any of mine. Verified $196/196$ in
consistent frames, $0$ failures, with all $12$ mismatched-frame cells and $4$
non-symmetric controls failing as predicted. Grade: \texttt{proved}.

\section{Status and grades}
\label{sec:grades}

Taken from \texttt{proofs/registry/conjecture-P.json}. Nothing here claims more than the
registry says. Grades are ordered
$\texttt{hunch} < \texttt{sketched} < \texttt{computed} < \texttt{checked-sober} <
\texttt{proved} < \texttt{lean-verified}$, plus \texttt{dead-end} and \texttt{in-progress}.

\begin{center}
\small
\sloppy\begin{longtable}{@{}p{0.53\textwidth}|p{0.20\textwidth}@{}}
\toprule
\textbf{Claim} & \textbf{Grade}\\
\midrule
\endhead
$\Psi^+(s_\mu)=\mathfrak s_\mu$ (Day 149 Thm A); $P_b=\sum_\mu K_{\mu'(2^b)}\mathfrak s_\mu$; $\tau$ is multiplication by $e_3$ & \texttt{proved}\\
Floor lemma: $V\mid \mathcal T(fV)$ with symmetric quotient, $\Psi(1)=1$ & \texttt{proved}\\
$\deg_u[T^n]H\le n$, i.e.\ (H2) in strong form & \texttt{proved}\\
Pieri-flow reformulation of (H1); integer-valued split-point partial (H1) & \texttt{proved}\\
\midrule
$\psi$ algebraic of degree $5$; the closed form; $Q$ (\textbf{this document}) & \texttt{computed}\\
$E$-positivity of $\psi$ to $Y^{33}$, $364$ monomials, $0$ negatives & \texttt{computed}\\
Narayana top layer at $E_3=0$, $n\le33$ & \texttt{computed}\\
Minimal-$E_3$ period-$2$ edge families; weight-homogeneity; $E_3\mid\psi_k$ for $k\ge3$ & \texttt{computed}\\
Conjecture P empirics: $[T^n]H$ for $n\le16$, $P_b$ for $b\le20$, $0$ negatives & \texttt{computed}\\
$s^*_\mu$ frame check ($23/23$); Kerov filtration $=$ $u$-degree filtration & \texttt{computed}\\
\midrule
Layer-induction top-down proof of Conjecture P & \texttt{sketched}\\
Conjecture P (root) & \texttt{in-progress}\\
(H1): $H\in\Ztwo[E_1,E_2,E_3][[T]]$ & \texttt{in-progress}\\
Minimality half of P (minimum $n+1$ attained at $E_1^n$) & \texttt{in-progress}\\
Identify $\psi$ combinatorially / the slice sequence & \texttt{in-progress}\\
\midrule
Catalan prediction $[W^n]\psi=C_{n+1}$ & \texttt{dead-end} (refuted)\\
Kerov character polynomial bridge (T1, T2) & \texttt{dead-end} (refuted)\\
Conjecture C ($F_P$ $\ell$-integral at separable $E$) & \texttt{dead-end} (false)\\
$J$-fraction route to (H1)/P & \texttt{dead-end}\\
Integer-valued split-point route pushed to all primes & \texttt{dead-end} (provably optimal, hence stalls)\\
$S_3$-orbit and Pieri-operator refinements of $P_b$ & \texttt{dead-end}\\
\bottomrule
\end{longtable}
\end{center}

\sloppy
\noindent\textbf{The registry has no node for the closed form yet.} As of this writing the
node \texttt{psi-slice-unidentified} is still \texttt{in-progress}, and among its recorded
negatives is ``not algebraic with $\deg_f\le3$, $\deg_W\le4$'' --- which is consistent with,
not contradicted by, the answer $\deg_f=5$, $\deg_W=3$. The scan box was simply too small,
as its own explicit caveat said. That node should now be closed by
Theorem~\ref{thm:main}/\ref{thm:slice} at grade \texttt{computed}.

\section*{What is still open}

\begin{enumerate}
\item \textbf{$E$-positivity of $\psi$ as a theorem.} $364$ positive monomials, no
certificate. $Q$ is not manifestly positive, so a different certificate is needed. The
period-$2$ minimal-$E_3$ families are the most model-shaped structure available.
\item \textbf{Turn \S\ref{sec:derivation} into a theorem.} The $\ltop$ bookkeeping behind
$\log\ltop_0(H)=\partial\Xi$ and $\theta\Xi=\tfrac12(P-E_1)$ has to be written out. Until it
is, everything here is \texttt{computed}.
\item \textbf{Identify $1,2,5,34,334,3958,52599,\dots$} beyond its algebraic equation.
Not in OEIS.
\item \textbf{The second degree-$5$ factor} of the resultant is unexplained. Presumably the
conjugate branch; not chased.
\item \textbf{The descent from layer $d$ to layer $d+1$} in the layer-induction attack. Still
holes.
\end{enumerate}

\end{document}
